% answers_short.tex —— 2015 年·陕西卷·文史类 · 简短答案速查
% 由 tools/make_short_answers.py 自 answers.tex 的 \gksolhead/\gkanswer 提取，
% 并与 2015/2015_数学真题_陕西文_答案要点.md 逐题比对后生成，勿手改（改源 MD 后重跑）。

\begin{gknotice}
\textbf{使用说明}：本节只给最终答案便于快速核对，完整推导见解析版。本题库无卷内答案：全部为项目独立求解、经第二人盲解对照（D 级），非官方标准答案；第 22–24 题均给出解答，固定副本未见选做说明。
\end{gknotice}

\gkpart{选择题}
\begin{center}\begin{tabular}{ccccccccccccc}
\toprule
题号 & 1 & 2 & 3 & 4 & 5 & 6 & 7 & 8 & 9 & 10 & 11 & 12 \\
\midrule
答案 & A & C & B & C & D & A & D & B & B & C & D & C \\
\bottomrule
\end{tabular}\end{center}

\gkpart{填空题}
\begin{itemize}[leftmargin=2.6em]
  \item[\textbf{13.}] $5$。
  \item[\textbf{14.}] $8$。
  \item[\textbf{15.}] $y=-\dfrac{1}{\mathrm{e}}$。
  \item[\textbf{16.}] $1-\dfrac{1}{2}+\dfrac{1}{3}-\cdots+\dfrac{1}{2n-1}-\dfrac{1}{2n}=\dfrac{1}{n+1}+\dfrac{1}{n+2}+\cdots+\dfrac{1}{2n}$。
\end{itemize}

\gkpart{解答题}
\begin{itemize}[leftmargin=2.6em]
  \item[\textbf{17.}] (1) $A=\dfrac{\pi}{3}$；(2) $\triangle ABC$ 的面积为 $\dfrac{3\sqrt{3}}{2}$
  \item[\textbf{18.}] (1) 证明见解析；(2) $a=6$
  \item[\textbf{19.}] (1) $\dfrac{13}{15}$；(2) $\dfrac{7}{8}$
  \item[\textbf{20.}] (1) $\dfrac{x^{2}}{2}+y^{2}=1$；(2) 证明见解析
  \item[\textbf{21.}] (1) $f_{n}'(2)=(n-1)2^{n}+1$；(2) 证明见解析
  \item[\textbf{22.}] (1) 证明见解析；(2) $\odot O$ 的直径为 $3$
  \item[\textbf{23.}] (1) $x^{2}+(y-\sqrt{3})^{2}=3$；(2) $P(3,0)$
  \item[\textbf{24.}] (1) $a=-3$，$b=1$；(2) 最大值为 $4$
\end{itemize}
